% Metadane pracy magisterskiej — edytuj ten plik, aby dostosować stronę tytułową,
% nagłówki i metadane PDF.

\newcommand{\thesisTitle}{Multimodalny, wieloagentowy system hybrydowy do wspomagania decyzji klinicznych}
\newcommand{\thesisTitleShort}{MedVoicePL --- system CDSS}

\newcommand{\thesisAuthor}{Jakub Boczar}
\newcommand{\thesisStudentId}{s23584}

\newcommand{\thesisSupervisor}{Dr Mateusz Żbikowski}
\newcommand{\thesisReviewer}{[tytuł naukowy] Imię Nazwisko}  % TODO: uzupełnić dane recenzenta

\newcommand{\thesisFaculty}{Wydział Informatyki}
\newcommand{\thesisDepartment}{Katedra: Systemy Inteligentne i Data Science}
\newcommand{\thesisSpecialization}{Specjalizacja: Data Science}

\newcommand{\thesisDegree}{Praca magisterska napisana pod kierunkiem:}

\newcommand{\thesisYear}{2026}
\newcommand{\thesisCity}{Warszawa}

% Metadane PDF (widoczne w właściwościach pliku)
\hypersetup{
  pdftitle={\thesisTitle},
  pdfauthor={\thesisAuthor},
  pdfsubject={Praca magisterska --- PJATK},
  pdfkeywords={CDSS, NLP, STT, anonimizacja, RAG, MedVoicePL}
}
